\documentclass{article} %210 mm × 297 mm
\usepackage[utf8]{inputenc}
\usepackage[a4paper,left=1.5cm, right=1.5cm, top=2cm, bottom=2cm]{geometry}
\usepackage{graphicx}
%\usepackage{blindtext}
\usepackage{multirow}
\usepackage{mhchem}
\usepackage{url}
\usepackage{subfigure}
\usepackage{amsfonts,latexsym} % para tener disponibilidad de diversos simbolos
\usepackage{enumerate}
%\usepackage{booktabs}
\usepackage{float}
%\usepackage{threeparttable}
\usepackage{array,colortbl}
%\usepackage{ifpdf}
%\usepackage{rotating}
\usepackage{cite}
\usepackage{stfloats}
\usepackage{url}
\usepackage{listings}
\usepackage{fancyhdr}

\usepackage{color}
%\usepackage{hyperref}
%\usepackage{wrapfig}
\usepackage{array}
%\usepackage{multirow}
%\usepackage{adjustbox}
%\usepackage{nccmath}
\usepackage[dvipsnames]{xcolor}


\newcommand{\titulo}{Abgabe 5; MfN II.}
\newcommand{\nombre}{Liliana Sanfilippo}
\newcommand{\FCE}{\textbf{SoSe 2025}}

 

\usepackage{fancyhdr}
\pagestyle{fancy}
\fancyhead{}
\fancyfoot{}
\fancyhead[LO]{\small\nombre}
\fancyhead[RE]{\small\titulo}
\fancyhead[LO]{\small\FCE}
\fancyfoot[RO,LE] {\thepage}

\setcounter{page}{1} %Número de la página, la editorial lo cambiará

\newcommand{\txc}[1]{\textcolor{blue}{#1}}

\usepackage[onehalfspacing]{setspace}
\begin{document}


\begin{tabular}{p{12cm}}
{{\Large\bf\titulo}}\\
\vspace{0.2cm}\\
 {\large\nombre}\\
\end{tabular}
\vspace{0.2cm}\\
\section{Aufgabe 5.1}
\subsection*{(a)}
\[
B^3 = \begin{pmatrix}
2 & -3 & 1 & -1 \\
3 & -6 & 0 & -3 \\
-1 & 3 & 1 & 2 \\
-3 & 6 & 0 & 3
\end{pmatrix} \cdot \begin{pmatrix}
2 & -3 & 1 & -1 \\
3 & -6 & 0 & -3 \\
-1 & 3 & 1 & 2 \\
-3 & 6 & 0 & 3
\end{pmatrix} \cdot \begin{pmatrix}
2 & -3 & 1 & -1 \\
3 & -6 & 0 & -3 \\
-1 & 3 & 1 & 2 \\ 
-3 & 6 & 0 & 3
\end{pmatrix} = \begin{pmatrix}
-3 & 9 & 3 & 6 \\
-3 & 9 & 3 & 6 \\
0 & 0 & 0 & 0 \\
3 & -9 & -3 & -6
\end{pmatrix} \cdot \begin{pmatrix}
2 & -3 & 1 & -1 \\
3 & -6 & 0 & -3 \\
-1 & 3 & 1 & 2 \\
-3 & 6 & 0 & 3
\end{pmatrix} = 0
\]
\subsection*{(b)}
%ges. eine Matrix $J$ mit $B=S^{-1}JS, \quad S \in K^{4\times4}$  und $J$ in Jordan-Normalform\\ 
%\txc{Was hat eine Jordan Basis damit zu tun? Wenn ich für B eine Jordan-Basis finde, kann ich daraus S bilden. Die Basis besteht aus den Eigenvektoren und der Jordan-Kette\\}
Charakteristisches Polynom von $B$ mit Laplace nach der dritten Spalte: 
\begin{align*}
    \chi_B(\lambda)  = det (B - \lambda I) & = det \begin{pmatrix}
2 - \lambda& -3 & 1 & -1 \\
3 & -6- \lambda & 0 & -3 \\
-1 & 3 & 1- \lambda & 2 \\
-3 & 6 & 0 & 3- \lambda
\end{pmatrix} \\
& = \sum_{i=1}^4 a_{i3} \cdot (-1)^{i+3} \cdot det(B'_{i3}) \\
& = a_{13} \cdot (-1)^{i+3} \cdot det(B'_{i3}) +  \sum_{i=2}^4 a_{i3} \cdot (-1)^{1+3} \cdot det(B'_{13}) \\
& = 1 \cdot (-1)^{4} \cdot det(B'_{13}) +  \sum_{i=2}^4 a_{i3} \cdot (-1)^{i+3} \cdot det(B'_{i3}) \\
& = det(B'_{13}) +  \sum_{i=2}^4 a_{i3} \cdot (-1)^{i+3} \cdot det(B'_{i3})  \\
& = det(B'_{13}) +  0 + \sum_{i=3}^4 a_{i3} \cdot (-1)^{i+3} \cdot det(B'_{i3})  \\ 
& = det(B'_{13}) +  a_{33} \cdot (-1)^{3+3} \cdot det(B'_{33}) + \sum_{i=3}^4 a_{i3} \cdot (-1)^{i+3} \cdot det(B'_{i3})  \\
& = det(B'_{13}) + (1-\lambda)  \cdot det(B'_{33}) + \sum_{i=4}^4 a_{i3} \cdot (-1)^{i+3} \cdot det(B'_{i3})  \\
& = det(B'_{13}) + (1-\lambda)  \cdot det(B'_{33})
\end{align*}
\[B'_{13} = \begin{pmatrix}
3 & -6- \lambda  & -3 \\
-1 & 3  & 2 \\
-3 & 6 &  3- \lambda
\end{pmatrix}  \]
\begin{align*}
    det(B'_{13}) & = 3\cdot3\cdot(-\lambda+3)+(-\lambda-6)\cdot2\cdot(-3)+(-3)\cdot(-1)\cdot6-(-3)\cdot3\cdot(-3)-6\cdot2\cdot3-(-\lambda+3)\cdot(-1)\cdot(-\lambda-6) \\
    & = 9\cdot(-\lambda+3)+(-\lambda-6)\cdot(-6)+18-27-36-(-\lambda+3)\cdot(-1)\cdot(-\lambda-6) \\
    & = -9\lambda + 27 + 
    6 \lambda + 36 
    +18-27-36
   + \lambda^2 + 3 \lambda - 18 \\
    & = -9\lambda +
    6 \lambda
   + \lambda^2 + 3 \lambda \\
     & =  \lambda^2  \\
\end{align*}
Also 
\[
 \chi_B(\lambda)  = \lambda^2  + (1-\lambda)  \cdot det(B'_{33})
\]
\[
B'_{33} = \begin{pmatrix}
2 - \lambda& -3  & -1 \\
3 & -6- \lambda & -3 \\
-3 & 6 & 3- \lambda
\end{pmatrix} 
\]
\begin{align*}
    det(B'_{33}) & = (-\lambda+2)\cdot(-\lambda-6)\cdot(-\lambda+3)
    +(-3)\cdot(-3)\cdot(-3)
    +(-1)\cdot3\cdot6
    -(-3)\cdot(-\lambda-6)\cdot(-1) \\
    & -6\cdot(-3)\cdot(-\lambda+2)
    -(-\lambda+3)\cdot3\cdot(-3) \\
    & = -\lambda^3 - \lambda^2 + 24 \lambda - 36 
    - 27  -18 
    + 3 \lambda + 18
    + 36 - 18 \lambda
    + 27 - 9 \lambda \\
    & = -\lambda^3 - \lambda^2 + 24 \lambda 
    + 3 \lambda 
  -18  \lambda
    - 9 \lambda \\
      & = -\lambda^3 - \lambda^2
\end{align*}
\[
 \chi_B(\lambda)  = \lambda^2  + (1-\lambda)  \cdot ( -\lambda^3 - \lambda^2) = \lambda^2 +  \lambda^4 - \lambda^2 =  \lambda^4 
\]
Eigenwerte finden: \\
$\lambda = 0$ mit algebraischer Vielfachheit vier.  \\
Da der Eigenwert 0 ist, sind hier Kern und die Menge der Eigenwerte identisch. \\ 
Kerne berechnen:   Dafür setze 
\[
 0 = B v = \begin{pmatrix}
2 & -3 & 1 & -1 \\
3 & -6 & 0 & -3 \\
-1 & 3 & 1 & 2 \\
-3 & 6 & 0 & 3
\end{pmatrix}  \cdot v
\]
B in Zeilenstufenform bringen. 
\[
\begin{pmatrix}
2 & -3 & 1 & -1 \\
3 & -6 & 0 & -3 \\
-1 & 3 & 1 & 2 \\
-3 & 6 & 0 & 3
\end{pmatrix} \overset{z4=z4+z2}{\rightsquigarrow}
\begin{pmatrix}
2 & -3 & 1 & -1 \\
3 & -6 & 0 & -3 \\
-1 & 3 & 1 & 2 \\
0 & 0 & 0 & 0
\end{pmatrix} \overset{z2=3*z3+z2}{\rightsquigarrow} \begin{pmatrix}
2 & -3 & 1 & -1 \\
0 & 3 & 3 & 3 \\
-1 & 3 & 1 & 2 \\
0 & 0 & 0 & 0
\end{pmatrix}  \overset{z3=0.5*z1+z2}{\rightsquigarrow} \begin{pmatrix}
2 & -3 & 1 & -1 \\
0 & 3 & 3 & 3 \\
0 & 1,5 & 1,5 & 1,5 \\
0 & 0 & 0 & 0
\end{pmatrix} 
\]
\[
 \overset{z3=(-1)*z2+z3}{\rightsquigarrow} \begin{pmatrix}
2 & -3 & 1 & -1 \\
0 & 3 & 3 & 3 \\
0 & 0 & 0 & 0 \\
0 & 0 & 0 & 0
\end{pmatrix} 
 \overset{z2=z2/3}{\rightsquigarrow} \begin{pmatrix}
2 & -3 & 1 & -1 \\
0 & 1 & 1 & 1 \\
0 & 0 & 0 & 0 \\
0 & 0 & 0 & 0
\end{pmatrix} 
 \overset{z1=z1/2}{\rightsquigarrow} \begin{pmatrix}
1 & -1,5 & 0,5 & -0,5 \\
0 & 1 & 1 & 1 \\
0 & 0 & 0 & 0 \\
0 & 0 & 0 & 0
\end{pmatrix}  \overset{z1=z2*(3/2)+z1}{\rightsquigarrow} \begin{pmatrix}
1 & 0 & 2 & 1 \\
0 & 1 & 1 & 1 \\
0 & 0 & 0 & 0 \\
0 & 0 & 0 & 0
\end{pmatrix} 
\]
Sete $B \cdot v = 0$.
Gleichungssystem lösen.
\[
\begin{array}{rrrrl}
a &  & +2c & +d &=0 \\
 & b & +c & +d &=0
\end{array} \quad \Rightarrow \quad b = -c-d, \quad a = -b-2c
\]
Kern hat Dimension 2: 
\[
v_1 = \begin{pmatrix}
-2 \\ -1 \\ 1 \\0 
\end{pmatrix}, \quad v_2 = \begin{pmatrix}
-1 \\ -1 \\ 0 \\ 1
\end{pmatrix}
\]
Kern von $B^2$ berechnen: Dafür setze 
\[
 0 = B^2 v = \begin{pmatrix}
-3 & 9 & 3 & 6 \\
-3 & 9 & 3 & 6 \\
0 & 0 & 0 & 0 \\
3 & -9 & -3 & -6
\end{pmatrix}  \cdot v
\]
In Zeilenstufenform bringen: 
\[
\begin{pmatrix}
-3 & 9 & 3 & 6 \\
-3 & 9 & 3 & 6 \\
0 & 0 & 0 & 0 \\
3 & -9 & -3 & -6
\end{pmatrix}  \overset{z2=z2-z1}{\rightsquigarrow}
\begin{pmatrix}
-3 & 9 & 3 & 6 \\
0 & 0 & 0 & 0\\
0 & 0 & 0 & 0 \\
3 & -9 & -3 & -6
\end{pmatrix} \overset{z4=z4+z1}{\rightsquigarrow}
\begin{pmatrix}
-3 & 9 & 3 & 6 \\
0 & 0 & 0 & 0\\
0 & 0 & 0 & 0 \\
0 & 0 & 0 & 0
\end{pmatrix}
\]
Gleichung lösen. 
\begin{align*}
    0 &= -3a + 9b + 3c + 6d & :3 \\
    & = -a +3b +c+2d \\
    \Rightarrow a &= 3b +c+2d
\end{align*}
also 
\[
\begin{pmatrix}
3b +c+2d \\ b \\c \\d
\end{pmatrix}
\]
Also ist der Kern: 
\[
\begin{pmatrix}
3\\1\\0\\0
\end{pmatrix}, \quad \begin{pmatrix}
1\\0\\1\\0
\end{pmatrix}, \quad \begin{pmatrix}
2\\0\\0\\1
\end{pmatrix}
\]
Und hat Dimension 3. \\
Da $A^3 = 0$ sind ab dem exponenten die Dimensionen des Kerns immer 4: \[
\begin{pmatrix}
1\\0\\0\\0
\end{pmatrix}, 
\quad \begin{pmatrix}
0\\1\\0\\0
\end{pmatrix},\quad \begin{pmatrix}
0\\0\\1\\0
\end{pmatrix}, \quad \begin{pmatrix}
0\\0\\0\\1
\end{pmatrix}
\]
Wenn man die Kerne von $A, A^2 und A^3$ berechnet, bekommt man heraus, dass 2, 3 und 4 sind. Also muss es einen Jordan-Block der Größe 3 geben und zwei Blöcke insgesamt. \\
Also: 
\[
J = \begin{pmatrix}
0&1&0&0\\
0&0&1&0\\
0&0&0&0\\
0&0&0&0
\end{pmatrix}
\]

\subsection*{(c)}
Zuerst Jordanbasis finden. \\
Dafür wählen wir einen Eigenvektor aus dem Kern mit der höchsten Dimension (einen veragllgemeinerten Eigenvektor).  \\
Wir wählen $v_4 = \begin{pmatrix}
1\\0\\0\\0
\end{pmatrix}$, da das (u.a.) in keinem der anderen Kerne enthalten ist. Dann setzen wir $v_3 = B \cdot v_4, \quad v_2 = B^2 \cdot v_4$: 
\[
v_3 = B \cdot v_4 : \begin{pmatrix}
2 & -3 & 1 & -1 \\
3 & -6 & 0 & -3 \\
-1 & 3 & 1 & 2 \\
-3 & 6 & 0 & 3
\end{pmatrix} \cdot  \begin{pmatrix}
1\\0\\0\\0
\end{pmatrix}  = \begin{pmatrix}
 2\\3\\-1\\-3
\end{pmatrix}
\]
\[
v_2 = B^2 \cdot v_4 : \begin{pmatrix}
-3 & 9 & 3 & 6 \\
-3 & 9 & 3 & 6 \\
0 & 0 & 0 & 0 \\
3 & -9 & -3 & -6
\end{pmatrix} \cdot  \begin{pmatrix}
 2\\3\\-1\\-3
\end{pmatrix}  = \begin{pmatrix}
-3\\-3\\0\\-3
\end{pmatrix}
\]
$v_1$ wählen wir aus dem Kern von B, hier $v_1 = \begin{pmatrix}
-2 \\ -1 \\ 1 \\0 
\end{pmatrix}$. \\
Dann ist die Jordan Basis: 
\[
\begin{pmatrix}
-2 \\ -1 \\ 1 \\0 
\end{pmatrix},
\begin{pmatrix}
1\\0\\0\\0
\end{pmatrix}, \begin{pmatrix}
 2\\3\\-1\\-3
\end{pmatrix}, 
\begin{pmatrix}
-3\\-3\\0\\-3
\end{pmatrix}
\]
Also
\[
S = \begin{pmatrix}
-3 & 2 & 1 & -2 \\
-3 & 3 & 0 & -1 \\
0 & -1 & 0 & 1 \\
3 & -3 & 0 & 0
\end{pmatrix}
\]



\section{Aufgabe 5.2}
Wenn das charakteristische Polynom $x^4$ ist, hat man den Eigenwert $0$ mit algebraischer Vielfachheit 4 (siehe vorherige Aufgabe). Das heißt die Dimensionen der Jordan-Blöcke müssen auf 4 summieren. 
Man kann 4 auf fünf verschiedene Arten und Weisen zerlegen: 
\begin{align*}
    4:& \begin{pmatrix}
    0&1&0&0\\
    0&0&1&0\\
    0&0&0&1\\
    0&0&0&0
    \end{pmatrix} &
    3+1:& \begin{pmatrix}
    0&1&0&0\\
    0&0&1&0\\
    0&0&0&0\\
    0&0&0&0
    \end{pmatrix} &
    2+2:& \begin{pmatrix}
    0&1&0&0\\
    0&0&0&0\\
    0&0&0&1\\
    0&0&0&0
    \end{pmatrix} \\
     2+1+1:& \begin{pmatrix}
    0&1&0&0\\
    0&0&0&0\\
    0&0&0&0\\
    0&0&0&0
    \end{pmatrix} &
     1+1+1+1:& \begin{pmatrix}
    0&1&0&0\\
    0&0&0&0\\
    0&0&0&0\\
    0&0&0&0
    \end{pmatrix}
\end{align*}

\section{Aufgabe 5.3}
\[
A = \begin{pmatrix}
3 & -1 & 4 & -1 \\
-4 & 0 & -4 & 1 \\
0 & 0 & -1 & 0 \\
5 & 5 & 5 & 4
\end{pmatrix}
\]
\subsection{(a)}
Charakteristisches Polynom von $A$ mit Laplace nach der dritten Zeile:
\begin{align*}
       \chi_A(\lambda)  = det (A - \lambda I) & = det \begin{pmatrix}
       3 - \lambda & -1 & 4 & -1 \\
-4 & - \lambda & -4 & 1 \\
0 & 0 & -1- \lambda & 0 \\
5 & 5 & 5 & 4- \lambda
       \end{pmatrix} \\
       & = \sum_{j=1}^4 a_{3j} \cdot (-1)^{3+j} \cdot det(A'_{3j}) \\
       & = a_{33} \cdot (-1)^{3+3} \cdot det(A'_{33})\\
       & = (-1- \lambda) \cdot det(A'_{33})
\end{align*}
\[
A'_{33} = \begin{pmatrix}
       3 - \lambda & -1 & -1 \\
-4 & - \lambda & 1 \\
5 & 5 & 4- \lambda
\end{pmatrix}
\]
\begin{align*}
    det(A'_{33}) & = (-\lambda+3)\cdot(-\lambda)\cdot(-\lambda+4)
    +(-1)\cdot1\cdot5
    +(-1)\cdot(-4)\cdot5
    -5\cdot(-\lambda)\cdot(-1)
    -5\cdot1\cdot(-\lambda+3)
    -(-\lambda+4)\cdot(-4)\cdot(-1) \\
    & = -\lambda^3 + 7 \lambda^2 - 12 \lambda
    -5
    +20
    -5\lambda
    +5 \lambda - 15
    + 4 \lambda - 16\\
     & = -\lambda^3 + 7 \lambda^2 - 8 \lambda - 16\\
\end{align*}
Also 
\[
  \chi_A(\lambda)  = (-1- \lambda) \cdot (-\lambda^3 + 7 \lambda^2 - 8 \lambda - 16) = (\lambda - 4)^2 (\lambda + 1)^2
  %= \lambda^4 - 6 \lambda^3 + \lambda^2 + 24 \lambda + 16
\]
Man kann $\lambda_{1} = -1$  und $\lambda_{2} = 4$ ablesen. Beide mit der algebraischen Vielfachheit 2. Also wird K die Form \[
\begin{pmatrix}
-1&?&0&0\\
0&-1&0&0\\
0&0&4&?\\
0&0&0&4
\end{pmatrix}
\]
haben. Kern von $A-\lambda I$ berechnen. \\
Für $\lambda = -1$:
\[
A +1 I = \begin{pmatrix}
4 & -1 & 4 & -1 \\
-4 & 1 & -4 & 1 \\
0 & 0 & 0 & 0 \\
5 & 5 & 5 & 5
\end{pmatrix} \overset{z2=z2+z1}{\rightsquigarrow} \begin{pmatrix}
4 & -1 & 4 & -1 \\
0 & 0 & 0 & 0 \\
0 & 0 & 0 & 0 \\
5 & 5 & 5 & 5
\end{pmatrix} \overset{z4 = z4/5}{\rightsquigarrow}
\begin{pmatrix}
4 & -1 & 4 & -1 \\
0 & 0 & 0 & 0 \\
0 & 0 & 0 & 0 \\
1 & 1 & 1 & 1
\end{pmatrix}
\]
\[
\overset{z1 = z1+z4}{\rightsquigarrow}
\begin{pmatrix}
5 & 0 & 5 & 0 \\
0 & 0 & 0 & 0 \\
0 & 0 & 0 & 0 \\
1 & 1 & 1 & 1
\end{pmatrix}
\overset{z1 = z1/5}{\rightsquigarrow}
\begin{pmatrix}
1 & 0 & 1 & 0 \\
0 & 0 & 0 & 0 \\
0 & 0 & 0 & 0 \\
1 & 1 & 1 & 1
\end{pmatrix}
\overset{z4 = z4-z1}{\rightsquigarrow}
\begin{pmatrix}
1 & 0 & 1 & 0 \\
0 & 0 & 0 & 0 \\
0 & 0 & 0 & 0 \\
0 & 1 & 0 & 1
\end{pmatrix}
\]
Setze  $(A−\lambda I) \cdot v = 0$ un dlöse das Gleichungssystem 
\[
\begin{array}{rrrrl}
 a &  & +c &  & = 0\\
& b & & +d & =0
\end{array} \Rightarrow a = -c, b =-d 
\]
Also ist der Kern Zwei dimensional: 
\[
\begin{pmatrix}
-1\\0\\1\\0
\end{pmatrix}, \begin{pmatrix}
0\\-1\\0\\1
\end{pmatrix}
\]
\[
(A +1 I)^2 = \begin{pmatrix}
1 & 0 & 1 & 0 \\
0 & 0 & 0 & 0 \\
0 & 0 & 0 & 0 \\
0 & 1 & 0 & 1
\end{pmatrix} = (A +1 I)^3 = (A +1 I)^4
\]
Also sind für $\lambda = -1$ die Jordan-Ketten maximal Eins lang und K muss die form haben 
\[
\begin{pmatrix}
-1&0&0&0\\
0&-1&0&0\\
0&0&4&?\\
0&0&0&4
\end{pmatrix}
\]
\[
A -4 I = \begin{pmatrix}
-1 & -1 & 4 & -1 \\
-4 & -4 & -4 & 1 \\
0 & 0 & -5 & 0 \\
5 & 5 & 5 & 0
\end{pmatrix} \overset{z4 = z4 + (z1+z2)}{\rightsquigarrow}
\begin{pmatrix}
-1 & -1 & 4 & -1 \\
-4 & -4 & -4 & 1 \\
0 & 0 & -5 & 0 \\
0 & 0 & 5 & 0
\end{pmatrix} \overset{z4=z4+z3}{\rightsquigarrow}
\begin{pmatrix}
-1 & -1 & 4 & -1 \\
-4 & -4 & -4 & 1 \\
0 & 0 & -5 & 0 \\
0 & 0 & 0 & 0
\end{pmatrix} 
\]
\[
 \overset{z3=z3/(-5)}{\rightsquigarrow}
\begin{pmatrix}
-1 & -1 & 4 & -1 \\
-4 & -4 & -4 & 1 \\
0 & 0 & 1 & 0 \\
0 & 0 & 0 & 0
\end{pmatrix} \overset{z2=(z2+z1)/(-5)}{\rightsquigarrow}
\begin{pmatrix}
-1 & -1 & 4 & -1 \\
1 & 1 & 0 & 0 \\
0 & 0 & 1 & 0 \\
0 & 0 & 0 & 0
\end{pmatrix} 
\overset{z1 = (z1+z2-4*z3)/(-1)}{\rightsquigarrow}
\begin{pmatrix}
0 & 0 & 0 & 1 \\
1 & 1 & 0 & 0 \\
0 & 0 & 1 & 0 \\
0 & 0 & 0 & 0
\end{pmatrix} 
\]
Gleichungssystem mit A * v = 0 lösen. 
\[
\Rightarrow c = d = 0, \quad a = -b
\]
Also ist der Kern eindimensional: 
\[
\begin{pmatrix}
-1\\1\\0\\0
\end{pmatrix}
\]
\[
(A -4 I )^2 =
\begin{pmatrix}
0 & 0 & 0 & 0 \\
1 & 1 & 0 & 1 \\
0 & 0 & 1 & 0 \\
0 & 0 & 0 & 0
\end{pmatrix} = (A -4 I )^3 
\]
Der Kern von $(A -4 I )^2$ (GS lösen):
\[
\begin{array}{rrrrl}
     a&+b&&+d &=0\\
     &&c&&=0
\end{array} \Rightarrow a = -b-d, \quad c = 0
\]
Kern ist zweidimensional mit
\[
\begin{pmatrix}
-1\\1\\0\\0
\end{pmatrix}, \begin{pmatrix}
-1\\0\\0\\1
\end{pmatrix}
\]
Die Dimensionen können nicht höher werden, also hat der größte Jordanblock die Größe 2 und K ist: 
\[
\begin{pmatrix}
-1&0&0&0\\
0&-1&0&0\\
0&0&4&1\\
0&0&0&4
\end{pmatrix}
\]






\end{document}